% ============================================================
%  US_TDA_clean.tex  -  AJUR submission version
%  Cleaned from Quarto-generated output.
%  Compile with XeLaTeX (required for EB Garamond via fontspec).
%
%  Figure files expected in:  figures/
%    fig-newGraph-1.pdf   fig-newGraph-2.pdf
%    fig-EDA-1.pdf        fig-Fisher-1.pdf
%    fig-PPI-CPI_Gap-1.pdf  fig-Wage_Growth-1.pdf
%
%  Recommended R export settings (in your .qmd chunks):
%    Single figures (fig-EDA, fig-Fisher, etc.):
%      fig-width: 6.73   fig-height: <match original>   dpi: 300
%    Side-by-side pair (fig-newGraph):
%      fig-width: 8.42   fig-height: <match original>   dpi: 300
% ============================================================

\documentclass[10pt, letterpaper]{article}

% ------ Page geometry (AJUR: 1" T/L/B, 0.5" R) ----------------
\usepackage[top=1in, left=1in, bottom=1in, right=0.5in]{geometry}

% ------ Font: EB Garamond (requires XeLaTeX or LuaLaTeX) -------
\usepackage{fontspec}
\setmainfont{EB Garamond}
\usepackage{anyfontsize}

% ------ Math ---------------------------------------------------
\usepackage{amsmath, amssymb}

% ------ Microtypography ----------------------------------------
\usepackage{microtype}

% ------ Spacing & paragraph style (AJUR: single, no indent) ----
\usepackage{setspace}
\setstretch{1}
\frenchspacing
\setlength{\parindent}{0pt}
\setlength{\parskip}{6pt plus 2pt minus 1pt}
\setlength{\emergencystretch}{3em}

% ------ Section headings: bold, ALL CAPS already in text, unnumbered
\setcounter{secnumdepth}{-\maxdimen}
\usepackage{titlesec}
\titleformat{\section}{\normalfont\bfseries}{}{0em}{}
\titlespacing{\section}{0pt}{12pt plus 4pt minus 2pt}{4pt plus 1pt minus 1pt}
\titleformat{\subsection}{\normalfont\itshape}{}{0em}{}
\titlespacing{\subsection}{0pt}{8pt}{2pt}

% ------ Captions: AJUR size-8, bold label, period separator ----
\usepackage{caption}
\usepackage{subcaption}
\captionsetup{
  font            = {small},
  labelfont       = {bf},
  labelsep        = period,
  justification   = raggedright,
  singlelinecheck = false
}

% ------ Tables -------------------------------------------------
\usepackage{longtable, booktabs, array, colortbl, multirow, calc}
\usepackage[table]{xcolor}
\usepackage{float}
\usepackage{etoolbox}
\makeatletter
\patchcmd\longtable{\par}{\if@noskipsec\mbox{}\fi\par}{}{}
\makeatother
\IfFileExists{footnotehyper.sty}{\usepackage{footnotehyper}}{\usepackage{footnote}}
\makesavenoteenv{longtable}

% ------ Graphics -----------------------------------------------
\usepackage{graphicx}
\makeatletter
\def\fps@figure{htbp}
\makeatother

% ------ Footnotes: bold uppercase letters (AJUR style) ---------
\renewcommand{\thefootnote}{\textbf{\Alph{footnote}}}

% ------ Hyperlinks ---------------------------------------------
\usepackage{bookmark}
\usepackage{xurl}
\urlstyle{same}
\usepackage[draft]{hyperref}
\hypersetup{
  colorlinks = true,
  linkcolor  = {blue},
  filecolor  = {Maroon},
  citecolor  = {Blue},
  urlcolor   = {Blue}
}

% ------ Tight list spacing (Pandoc helper) ---------------------
\providecommand{\tightlist}{%
  \setlength{\itemsep}{0pt}\setlength{\parskip}{0pt}}

% ------ Citation machinery (Pandoc citeproc) -------------------
\NewDocumentCommand\citeproctext{}{}
\NewDocumentCommand\citeproc{mm}{%
  \begingroup\def\citeproctext{#2}\cite{#1}\endgroup}
\makeatletter
  \let\@cite@ofmt\@firstofone
  \def\@biblabel#1{}
  \def\@cite#1#2{{#1\if@tempswa , #2\fi}}
\makeatother

\newlength{\cslhangindent}
\setlength{\cslhangindent}{1.5em}
\newlength{\csllabelwidth}
\setlength{\csllabelwidth}{3em}
\newenvironment{CSLReferences}[2]
  {\begin{list}{}{%
    \setlength{\itemindent}{0pt}
    \setlength{\leftmargin}{0pt}
    \setlength{\parsep}{0pt}
    \ifodd #1
      \setlength{\leftmargin}{\cslhangindent}
      \setlength{\itemindent}{-1\cslhangindent}
    \fi
    \setlength{\itemsep}{#2\baselineskip}}}
  {\end{list}}
\newcommand{\CSLBlock}[1]{\hfill\break\parbox[t]{\linewidth}{\strut\ignorespaces#1\strut}}
\newcommand{\CSLLeftMargin}[1]{\parbox[t]{\csllabelwidth}{\strut#1\strut}}
\newcommand{\CSLRightInline}[1]{\parbox[t]{\linewidth - \csllabelwidth}{\strut#1\strut}}
\newcommand{\CSLIndent}[1]{\hspace{\cslhangindent}#1}

\begin{document}


{
\centering
{\fontsize{16}{19}\selectfont\textbf{Mapping the Shape of the U.S. Economy: A Topological Data Analysis Approach with Ball Mapper}}\par
\vspace{6pt}
}

\emph{Ryan Johnson}\textsuperscript{*}

\emph{Graduate, Department of Mathematics, University of Alaska
Anchorage, Anchorage, AK}

Student: \emph{johnson.ryan-0@pm.me}\textsuperscript{*} \newline Mentor:
\emph{scook25@alaska.edu}

\section{ABSTRACT}\label{abstract}

Topological Data Analysis (TDA) is a developing data analysis method.
Currently, a large body of TDA research utilizes the traditional Mapper
algorithm. This paper aims to show an entry-level path into TDA, and to
expand the body of literature using Ball Mapper--an algorithm adjacent
to traditional Mapper. Using U.S. macroeconomic data from the Bureau of
Economic Analysis (BEA), the Federal Reserve Bank (FRB), and the Bureau
of Labor Statistics (BLS), we begin by showcasing Ball Mapper's use in
exploratory data analysis. We then show how we can effectively use Ball
Mapper in combination with other derived economic measures: the Real
Interest Rate, PPI-CPI Gap, and Real Wage Growth. Our results reveal an
underlying structure in the U.S. economy from 1960 to 2024. Illustrated
through Ball Mapper's graphs, we find that economic expansion contains
distinct stages, that peaks and troughs are structurally separate from
the broader expansion cycle, and that no two economic contractions are
alike. Notably, despite not being included as an input, our analyses
uncovered meaningful relationships between years, which suggests that
the relationships seen are inherent to the economic data itself.

\section{KEYWORDS}\label{keywords}

Topological Data Analysis; Topology; TDA; Ball Mapper; Mapper; Data
Science; Economics; Macroeconomics

\section{INTRODUCTION}\label{introduction}

Topological Data Analysis (TDA) is an emerging field of data analysis
that offers a unique vantage point for exploring complex data. Broadly,
traditional TDA applications use two tools: an algorithm called Mapper,
whose output is represented by a mathematical graph, and an analysis
technique called Persistent
Homology.\textsuperscript{\citeproc{ref-madukpe2025mapper}{\textbf{1}}}\footnote{A mathematical graph, from Graph Theory, is a diagram that consists of nodes and edges.}
These two tools are not necessarily dependent on one another and both
provide information about data in their own right. Moreover, Mapper and
Persistent Homology are central to TDA's core argument: data contains an
underlying shape, and this shape can provide us with qualitative, and
sometimes quantitative, insights about large multidimensional
datasets.\textsuperscript{\citeproc{ref-chazal2021introduction}{\textbf{2}}}

At a high level, Mapper is used to explore and visualize
high-dimensional data. Briefly, the Mapper algorithm achieves this by
producing a mathematical graph using a filter function. This graph
presents a picture analogous to taking an x-ray of the data, revealing
the skeleton of the dataset; it allows one to find structural
relationships that are not obvious at first glance. Persistent
Homology's goal is to uncover topological features of a dataset
(components, holes, voids, etc.) through various scales. As an example,
think about viewing a small picture made of dots as you move it from
touching your nose to arm's length away. Too close, you only see
individual dots; too far away, the dots merge together into a singular
shape. Persistent Homology takes observational notes about how the
picture changes across all the various distances.

Moving forward, we will use the terms ``traditional Mapper'' and
``Mapper'' interchangeably. When we use the term ``traditional Mapper,''
we are referring to the methodology for Mapper presented by Singh et
al.\textsuperscript{\citeproc{ref-singh2007topological}{\textbf{3}}} in
2007. This is an important distinction because the concept of the Mapper
methodology has been implemented in various ways by changing the filter
function and other parameters.

Although Persistent Homology provides insight into the topological
features our data contains, the focus of this paper is on the structural
exploration that Mapper and Mapper variations offer. Specifically, we
are examining various macroeconomic indicators of the United States
(U.S.) using a Mapper-adjacent algorithm called Ball Mapper (BM). One
advantage of BM is that it requires fewer parameters than traditional
Mapper, allowing for easier implementation.

In traditional Mapper, one must pass the data through three stages, all
with their own parameters, in order to generate a graph. Ball Mapper,
however, only needs one parameter before producing a
graph.\textsuperscript{\citeproc{ref-dlotko2019ball}{\textbf{4}}} An
important trade-off between Ball Mapper and traditional Mapper is that
while BM reduces the number of steps, one loses the ability to fine-tune
outputs. Additionally, traditional Mapper literature presents itself as
requiring a mathematical background, such as an introductory course in
Topology, or conceptual knowledge of data science methods and
algorithms. Although it is not impossible to learn TDA without the above
backgrounds, TDA might seem unnecessarily complicated for data analysis,
or impossible to decipher.

One major motivation for this paper came from a paper by D\l{}otko et
al.\textsuperscript{\citeproc{ref-dlotko2019mapping}{\textbf{5}}} The
paper was a two-pronged macroeconomic analysis using Ball Mapper. The
first prong used a dataset of five macroeconomic variables from 16
countries from the 1870s to 2017. The authors were interested in
comparing how countries have evolved over time, their transformation
from the Great Depression Era, and various views on wealth and
inequality. The second prong was to look at the relationship between
private credit growth and the GDP of various countries. This paper was
the only one in the literature that we could find that combines both
macroeconomics and Ball Mapper.

This analysis will expand the small body of literature whose main focus
is applications with Ball
Mapper.\textsuperscript{\citeproc{ref-madukpe2025mapper}{\textbf{1}}} We
also hope that it will serve as an on-ramp for anyone interested in TDA
who feels inundated with jargon during early-stage research.

Using U.S. macroeconomic data from the Bureau of Economic Analysis
(BEA), the Federal Reserve Bank (FRB), and the Bureau of Labor
Statistics (BLS), we plan to achieve these objectives in two stages. The
first stage is to show Ball Mapper's use in exploratory data analysis
(EDA) by looking at structural observations, coloration, and the size
variation of nodes in our graph. Then, using what we demonstrated in our
first stage, we extend that framework to other derived macroeconomic
measures: the Real Interest
Rate,\footnote{Approximation using the Fisher Equation.}\textsuperscript{\citeproc{ref-mcclung2024fisher}{\textbf{6}},
\citeproc{ref-michaelides2024fisher}{\textbf{7}}} PPI-CPI
Gap,\footnote{Producer Price Index: Finished Goods less Consumer Price Index.}\textsuperscript{\citeproc{ref-bls_ppi_cpi_comparison}{\textbf{8}}}
and Real Wage
Growth.\textsuperscript{\citeproc{ref-sullivan1997trends}{\textbf{9}}}\footnote{Compensation less Inflation (CPI).}

\section{METHODS AND PROCEDURES}\label{methods-and-procedures}

\emph{Data Selection \& Preparation}

We pulled data from three publicly available U.S. government sources:
The Bureau of Economic Analysis (BEA), The Federal Reserve Bank (FRB),
and The Bureau of Labor Statistics
(BLS).\textsuperscript{\citeproc{ref-bea_data}{\textbf{10}},
\citeproc{ref-fred_data}{\textbf{11}}} We gathered the data using R via
two application programming interfaces (APIs): one for data from the
BEA, and the other from the Federal Reserve Economic Data (FRED)
API.\textsuperscript{\citeproc{ref-bea_r_package}{\textbf{12}},
\citeproc{ref-fredr_package}{\textbf{13}}}\footnote{FRED aggregates data from national and international sources, as well as public and private sources.}
We additionally used recession dates based on business cycle
contractions and expansions provided by the National Bureau of Economic
Research (NBER). These agencies were selected because they are
authoritative sources for U.S. economic data. Their widespread use in
both the private and public sectors gives us high confidence in the
accuracy and integrity of the
data.\textsuperscript{\citeproc{ref-hughes2019value}{\textbf{14}}}

\begin{table}[H]

\centering{

\fontsize{8.0pt}{10.0pt}\selectfont
\begin{tabular*}{1\linewidth}{@{\extracolsep{\fill}}lcccc>{\centering\arraybackslash}p{\dimexpr 90.00pt -2\tabcolsep-1.5\arrayrulewidth}c}
\toprule
Data Series & Series Abbreviation & Year Start & Year End & Frequency & Series & API \\ 
\midrule\addlinespace[2.5pt]
Gross Domestic Product & GDP & 1930 & 2024 & Annual & Table 1.1.1 & BEA \\ 
{\cellcolor[HTML]{F2F2F2}{Personal Income \& Its Disposition}} & {\cellcolor[HTML]{F2F2F2}{PID}} & {\cellcolor[HTML]{F2F2F2}{1948}} & {\cellcolor[HTML]{F2F2F2}{2024}} & {\cellcolor[HTML]{F2F2F2}{Annual}} & {\cellcolor[HTML]{F2F2F2}{Table 2.1}} & {\cellcolor[HTML]{F2F2F2}{BEA}} \\ 
Foreign Direct Investment & FDI & 2014 & 2024 & Annual & Direct Investment and MNE & BEA \\ 
{\cellcolor[HTML]{F2F2F2}{Federal Funds Rate}} & {\cellcolor[HTML]{F2F2F2}{FFR}} & {\cellcolor[HTML]{F2F2F2}{1955}} & {\cellcolor[HTML]{F2F2F2}{2025}} & {\cellcolor[HTML]{F2F2F2}{Annual}} & {\cellcolor[HTML]{F2F2F2}{RIFSPFFNA}} & {\cellcolor[HTML]{F2F2F2}{FRED}} \\ 
Employment Cost Index & ECI & 2001 & 2025 & Quarterly & ECIALLCIV & FRED \\ 
{\cellcolor[HTML]{F2F2F2}{Consumer Price Index}} & {\cellcolor[HTML]{F2F2F2}{CPI}} & {\cellcolor[HTML]{F2F2F2}{1947}} & {\cellcolor[HTML]{F2F2F2}{2025}} & {\cellcolor[HTML]{F2F2F2}{Monthly}} & {\cellcolor[HTML]{F2F2F2}{CPIAUCSL}} & {\cellcolor[HTML]{F2F2F2}{FRED}} \\ 
New Privately-Owned Housing Units Started & Housing Starts & 1959 & 2025 & Monthly & HOUST & FRED \\ 
{\cellcolor[HTML]{F2F2F2}{Producer Price Index - Finished Goods}} & {\cellcolor[HTML]{F2F2F2}{PPI}} & {\cellcolor[HTML]{F2F2F2}{1947}} & {\cellcolor[HTML]{F2F2F2}{2025}} & {\cellcolor[HTML]{F2F2F2}{Monthly}} & {\cellcolor[HTML]{F2F2F2}{WPSFD49207}} & {\cellcolor[HTML]{F2F2F2}{FRED}} \\ 
Unemployment Rate & Unrate & 1948 & 2025 & Monthly & UNRATE & FRED \\ 
\bottomrule
\end{tabular*}

}

\caption{\label{tbl-Data_Sources}Summary of data series, time period,
and sources for this analysis.}

\end{table}%

Table~\ref{tbl-Data_Sources} shows all the data series we considered for
this analysis. We excluded Employment Cost Index (ECI) and Foreign
Direct Investment (FDI) due to their limited availability of years. Of
the remaining data series, New Privately-Owned Housing Units Started
(Housing Starts) had the smallest range and so set our lower bound for
the range years used (1960-2024). Our upper bound was the latest full
calendar year that was available. We wanted to represent the economy
from various vantage points, i.e., the broad economy, policy makers,
consumers, business owners, etc. Table~\ref{tbl-Analysis_Data} shows the
data series that we felt were representative of the economy as a whole
and had data for our time period of interest. Within each of these
series we used a subset of columns to construct our final dataset for
analysis. To see what parts of the economy we determined each series
plays, the column Functional Description summarizes them and can be
found in Table~\ref{tbl-Analysis_Data}.

For this analysis we are focusing on an annual time frame and data
transformation was needed. Looking at Table~\ref{tbl-Data_Sources}, not
all our data was provided on an annual basis.
Table~\ref{tbl-Analysis_Data} shows a summary of the transformation
taken and the details for each follow.

The first transformation we did was on Personal Income and Its
Disposition (subsequently, Personal Income). This data is reported in
nominal dollars and does not account for inflation. So, to account for
inflation we used the Consumer Price Index (CPI) to transform our values
from nominal dollars to real dollars
(Equation~\ref{eq-CPI_Real}).\textsuperscript{\citeproc{ref-dallasfed_nominal}{\textbf{15}}}
We then took the log difference (Equation~\ref{eq-deltaLog}) from the
output from Equation~\ref{eq-CPI_Real} to get an approximate percentage
change from the preceding year.

\begin{equation}\phantomsection\label{eq-CPI_Real}{
\text{Real Dollars} = \frac{\text{Nominal Dollars}}{\text{CPI}} \times 100
}\end{equation}

Calculating the year-over-year percent change using log differences
provides us with a few advantages. First, taking the log stabilizes the
variance and makes exponential trends across the series more linear.
Second, the difference between log values approximates the true
percentage change well for small to moderate values. Further, the log
difference allows for symmetric percent changes: increases or decreases
from year to year are of equal
magnitude.\textsuperscript{\citeproc{ref-hamilton2014logarithms}{\textbf{16}}}
An interesting question that has not been explored to our knowledge is
the impact on linearized data and Ball Mapper's distance metric, the
Euclidean distance. Further exploration of this question will be touched
on in the \emph{Discussion} section.

\begin{equation}\phantomsection\label{eq-deltaLog}{
\Delta\ln{(\text{Level})} = 
[\ln{(\text{Level}_{t})} - \ln{(\text{Level}_{t-1})}] \times 100
}\end{equation}

Our next transformations were applied to housing starts, Producer Price
Index - Finished Goods (PPI), and CPI. All three of these sources were
only provided on a monthly time frame, so to get an annual value for
them we used the simple arithmetic mean. Once these data were in annual
form, we then found the percent change from the previous year using the
log difference (Equation~\ref{eq-deltaLog}).

Housing starts data is given in a seasonally adjusted annual rate
(SAAR). To get this number the U.S. Census first takes the raw monthly
number of housing starts reported to their survey. Then they adjust that
number to account for seasonal patterns (e.g.~summertime booms or winter
lulls), and finally they multiply it by 12. How we interpret this final
number is: for the reporting month, it is the number of houses that
would be built at the end of a 12-month
period.\textsuperscript{\citeproc{ref-census_soc_methodology}{\textbf{17}}}

\begin{table}[H]

\centering{

\fontsize{8.0pt}{10.0pt}\selectfont
\begin{tabular*}{1\linewidth}{@{\extracolsep{\fill}}llll}
\toprule
Series Abbreviation & Economic Role & Functional Description & Transformation Applied \\ 
\midrule\addlinespace[2.5pt]
GDP & Growth Composition & Consumption, Investment, Gov., and Trade & None (Source in \% Change) \\ 
{\cellcolor[HTML]{F2F2F2}{PID}} & {\cellcolor[HTML]{F2F2F2}{Income Structure}} & {\cellcolor[HTML]{F2F2F2}{Wages, Entrepreneurship, Gov. Transfers}} & {\cellcolor[HTML]{F2F2F2}{Real Adjustment \& Log Difference}} \\ 
UnRate & Labor Market Dynamics & Labor Distress Level and Trend & None (Annual Rate) \& Simple Difference \\ 
{\cellcolor[HTML]{F2F2F2}{FFR}} & {\cellcolor[HTML]{F2F2F2}{Monetary Policy}} & {\cellcolor[HTML]{F2F2F2}{Cost of Capital and Borrowing Conditions}} & {\cellcolor[HTML]{F2F2F2}{None (Annual Rate) \& Simple Difference}} \\ 
PPI & Supply-Side Signal & Producer Costs \& Measures Cost-Push Inflation & 12-month Avg. \& Log Difference \\ 
{\cellcolor[HTML]{F2F2F2}{CPI}} & {\cellcolor[HTML]{F2F2F2}{Demand-Side Signal}} & {\cellcolor[HTML]{F2F2F2}{Cost of Living \& Measures Demand-Pull Inflation}} & {\cellcolor[HTML]{F2F2F2}{12-month Avg. \& Log Difference}} \\ 
Housing Starts & Leading Indicator & Physical Residential Production & 12-month Avg. \& Log Difference \\ 
\bottomrule
\end{tabular*}

}

\caption{\label{tbl-Analysis_Data}All years for analysis are 1960-2024.}

\end{table}%

Both PPI and CPI are economic indexes that represent a weighted average
of goods they respectively track. Although PPI has many different
series, the one used for this analysis tracks the price the producer
receives from the physical goods that it sells to businesses or
consumers. An important distinction to note about the PPI used here is
that the index captures goods which are finished or final in the
economic sense. Colloquially, it is similar to the wholesale price a
retailer might pay to a producer before marking up for retail sales.
However, it also includes goods which are sold to businesses, think
fully assembled computers or fleet
vehicles.\textsuperscript{\citeproc{ref-bls_ppi_handbook}{\textbf{18}}}
CPI on the other hand captures the consumer side and the retail price.
That is, the cost of the good after profit margin, transportation costs,
and other intermediary costs are calculated into a goods price sold at
the
register.\textsuperscript{\citeproc{ref-bls_cpi_concepts}{\textbf{19}}}
Broadly, both have been associated with measuring inflation for
producers and consumers. The more nuanced understanding is that index
changes can reflect many things: cost-push inflation, demand-pull
inflation, supply shocks, monetary policy inflation, and wage-price
spirals are just a few examples. It is our hope that Ball Mapper might
show the distinctive separation in the various aforementioned reasons
for the PPI and CPI indexes changing from year to year.

Finally, after all the transformations are complete, we used a total of
16 variables from our selected data series
(Table~\ref{tbl-Analysis_Data}) for our analysis. Again, we aimed to use
a comprehensive set of variables to represent the U.S. economy. A list
of variables used can be found in Table~\ref{tbl-Analysis_Dimensions} in
no particular order.

\begin{table}[H]

\centering{

\fontsize{8.0pt}{10.0pt}\selectfont
\begin{tabular*}{1\linewidth}{@{\extracolsep{\fill}}llll}
\toprule
 &  &  &  \\ 
\midrule\addlinespace[2.5pt]
PPI Finished Change & Personal Income & Compensation & Entrepreneurship \\ 
{\cellcolor[HTML]{F2F2F2}{Transfers}} & {\cellcolor[HTML]{F2F2F2}{Consumption}} & {\cellcolor[HTML]{F2F2F2}{Domestic Investment}} & {\cellcolor[HTML]{F2F2F2}{Government Spending}} \\ 
Exports & Imports & Inflation & Unemployment \\ 
{\cellcolor[HTML]{F2F2F2}{Unemployment Change}} & {\cellcolor[HTML]{F2F2F2}{Fed Rate}} & {\cellcolor[HTML]{F2F2F2}{Fed Rate Change}} & {\cellcolor[HTML]{F2F2F2}{Housing Change}} \\ 
\bottomrule
\end{tabular*}

}

\caption{\label{tbl-Analysis_Dimensions}Variables inculde in Ball Mapper
Point Cloud}

\end{table}%

\emph{Ball Mapper}

In-depth descriptions of Ball Mapper's theory have been covered in
various
papers.\textsuperscript{\citeproc{ref-dlotko2019ball}{\textbf{4}},
\citeproc{ref-rudkin2025ballmapper}{\textbf{20}},
\citeproc{ref-qiu2020refining}{\textbf{21}}} As the aim of our paper is
to provide an introductory foray into Ball Mapper and TDA in general, we
will omit some of the more technical details below.

Theoretically, Ball Mapper supposes we are given a dataset \(X\) with
\(K\) dimensions, \(N\) observations and some metric \(d\). Then to
create a Ball Mapper graph, our goal is to construct a collection of
balls, let's call it \(B\), with a constant radius, \(\epsilon\), such
that every observation of our dataset is covered by at least one ball in
\(B\). Mathematically, this is written as follows: Let \((X, d)\) be a
metric space where \(X \subset \mathbb{R}^K\) is our dataset and \(d\)
is a distance metric on \(X\). For a fixed radius \(\epsilon >0\), let
\(C \subseteq X\) be a set of center points. Then for each \(c \in C\),
we define a ball
\(b(c, \epsilon) = \{x \in X : d(c, x) \leq \epsilon\}\). Our collection
then is \(B = \{b(c, \epsilon) : c \in C\}\) and
\(\bigcup_{b \in B} b = X\).

In practice, our dataset \(X\) is called a Point Cloud. Our \(N\)
observations or rows are normally referenced as points, and our \(K\)
dimensions, the columns of \(X\), are denoted variables. Since Ball
Mapper is known to handle large multidimensional datasets well, it is
common to have point clouds with dimensions \(K > 2\). Ball Mapper was
originally written in R, but later on was adapted for use in Python. For
this analysis we used R and the Ball Mapper package aptly named
``BallMapper''.\textsuperscript{\citeproc{ref-ballmapper_package}{\textbf{22}}}

To construct the Ball Mapper (BM) graph, the balls of our graph are
created in an iterative manner and not a predetermined manner like the
theory defines. If we pre-defined a subset of ball centers, it would
take away from the driving idea behind TDA--that data has shape--as well
as start moving into the territory of other clustering methods such as
K-Means Clustering, where you select a number of clusters based on
statistical properties of the data. This iterative process allows our
point cloud to cluster in a ``natural'' way based on our choice of
epsilon and distance metric.

\begin{equation}\phantomsection\label{eq-euclideanDistance}{
d(x_i, x_j) = 
\sqrt{(x_{i_1} - x_{j_1})^2 + (x_{i_2} - x_{j_2})^2 + \ldots + (x_{i_k} - x_{j_k})^2} \leq \epsilon
}\end{equation}

The iterative ball creation is carried out in a linear fashion from the
first row to the last of the point cloud. The result of this linear
selection process is a greedy, sequential creation of ball centers. How
this works is that it takes the first point, let's call it \(c_1\), of
the point cloud and assigns it as the center of the first node in our
graph. Then using the default distance metric, the Euclidean Distance
(Equation~\ref{eq-euclideanDistance}), it compares every point in the
point cloud to \(c_1\) to create
\(b_1 = b(c_1, \epsilon) = \{x \in X : d(c_1, x) \leq \epsilon \}\).
After this first ball is created, Ball Mapper looks for the next row in
the point cloud that is not a member of the first ball, and assigns it
as the next center. This process continues until all the points are in a
ball and the union of balls covers our whole point cloud.

Hopefully, it is clear now why this process can be described as greedy
and sequential. It relies a great deal on the point cloud order as well
as the epsilon. In fact, the Ball Mapper graph changes for epsilons
sufficiently large enough to change at least one point's inclusion or
exclusion to a ball; the graph can also change for unique orders of the
point cloud rows. We touch on this in greater detail during our
\emph{Discussion}.

The last step to the Ball Mapper graph is to draw the edges if they
exist. The way the BM algorithm is structured, it allows for the
possibility of points to be in multiple balls. When two balls (also
called nodes) contain the same point, they are said to intersect. If two
nodes share at least one point, then an edge is created. If there is
more than one shared point, the number of shared points dictates the
edge strength. So, once we have all the nodes, edges, and edge
strengths, we then can construct a Ball Mapper graph. The Ball Mapper
algorithm described above can be broken down as follows:

\begin{enumerate}
\def\labelenumi{\arabic{enumi}.}
\tightlist
\item
  Select a point \(x_i \in X\) and assign it as \(c_i\).
\item
  Given \(\epsilon > 0\), construct a ball \(b_i(c_i, \epsilon)\) with a
  center \(c_i\) and radius \(\epsilon\) by checking if \(x_j\)
  satisfies \(d(c_i, x_j) \leq \epsilon : x_j \in X\) and \(d\) is the
  distance metric (Equation~\ref{eq-euclideanDistance}).
\item
  Place \(b_i\) in a set \(B\).
\item
  Select another \(x_n\) and repeat steps 1--3 such that
  \(x_n \neq x_j : x_j \in b_i\).
\item
  Draw an edge between \(b_i, b_j \in B : b_i \neq b_j\) if
  \(b_i \cap b_j \neq \emptyset\), weighting the edge based on the size
  of the intersection.
\end{enumerate}

\emph{Implementation}

To use Ball Mapper one only needs three parts: a point cloud, a variable
to color by, and an epsilon. For this analysis, our point cloud was made
up of 16 variables and 65 points, each representing a year from
1960--2024. Two of the variables, \emph{Unemployment Change} and
\emph{Fed Rate Change}, were created before any transformation took
place. We elected to create these because the Unemployment Rate
(Unemployment) and the Federal Funds Rate (Fed Rate) were the only two
dimensions which did not have a transformation and thus did not have a
year-over-year change like our other variables. Having the rate of
change for both Unemployment and the Fed Rate is also important because
it gives us additional information beyond the intrinsic value of each of
the measures. Both measures on their own give us a current look at the
economy: the Unemployment Rate telling us how many who are not working
but want
to,\textsuperscript{\citeproc{ref-bls_cps_concepts}{\textbf{23}}} and
the Fed Rate signaling the structural cost of
capital.\textsuperscript{\citeproc{ref-fed_monetary_goals}{\textbf{24}}}

There seems to be no generally accepted way of choosing epsilon in the
literature. Since Ball Mapper constructs maps by selecting a point and
creating a ball with a radius of epsilon, it is possible to select an
epsilon small enough to produce a graph where every point is a node of
only itself. Conversely, this should lead to an epsilon big enough such
that all of our points are included in one node. For our analysis here,
this was indeed the case. So to find an appropriate epsilon we first
started by finding lower and upper bounds.

What we look for in our lower bound is a graph with an epsilon small
enough to make every point in our point cloud a node of one point, a
singleton node, not connected to any point. For our upper bound we
similarly look for an epsilon so large that it creates a singular node
that houses all of our points. Additionally, to reduce some of the
computational time one can look for a lower bound where a connected
component starts to form, and an upper bound where we have a connected
component of only two nodes. This cuts off tail end graphs that
generally don't provide much substantial insight into the data.

For our analysis we found an initial lower and upper bound represented
by the interval \([0.38, 0.90]\). We then wrote a function to generate
around 100 graphs, reviewed them, and narrowed the upper and lower
bounds to \([0.40, 0.70]\). During this narrowing process we were
looking for interesting features: connected components forming or
dissolving, flares coming off of any components, and notable sizing or
coloration patterns. The resulting value of this process was \(0.474\).
We thought it contained all the aforementioned features. This final
epsilon is used for all the graphs generated in the analysis.

\begin{figure}[H]

\begin{minipage}{0.50\linewidth}

\centering{

\includegraphics[width=1\linewidth,height=\textheight,keepaspectratio]{figures/fig-newGraph-1.pdf}

}

\subcaption{\label{fig-newGraph-1}Ball Mapper graph using standard
function, ColorIgraphPlot.}

\end{minipage}%
%
\begin{minipage}{0.50\linewidth}

\centering{

\includegraphics[width=1\linewidth,height=\textheight,keepaspectratio]{figures/fig-newGraph-2.pdf}

}

\subcaption{\label{fig-newGraph-2}Ball Mapper graph using new graph
output created with ggraph.}

\end{minipage}%

\caption{\label{fig-newGraph}Ball Mapper graph colored by Year; epsilon
set to 0.511.}

\end{figure}%

Before describing some visual changes we made to the Ball Mapper output
graph, it is important to note that Ball Mapper will output the same
graph unless at least one of two things changes: the size of the point
cloud changes or the order of the point cloud. The former intuitively
makes sense. The algorithm is creating epsilon radius balls so adding or
subtracting any data will yield a different amount of points in at least
one ball. Depending on the data being changed or number of points, the
structure of the output might not change much. As an example, when we
were deciding on which PPI series to use, we had a multitude of choices.
Initially we used PPI - All Commodities but when reading the BLS
handbook on PPI, it specifically mentions that PPI - All Commodities is
susceptible to double counting. To avoid this issue we used PPI -
Finished Goods
instead.\textsuperscript{\citeproc{ref-bls_ppi_handbook}{\textbf{18}}}

The standard Ball Mapper package in R produces a sufficient graph for
analysis. The output object of the Ball Mapper algorithm is a series of
nodes, edges, edge weights, and other graph properties created by the
package
\emph{igraph}.\textsuperscript{\citeproc{ref-csardi2006igraph}{\textbf{25}}}
Separate from the igraph object, a graphing function that uses R's base
plot function and default coloring scheme produces the visual graph.
This is an advantage because when packages are written in Base R they
are usable to anyone who can install R. We, however, felt that this
graph output could be improved upon and recognized that we could use
alternative, and more comprehensive packages to create a graph,
\emph{ggraph}.\textsuperscript{\citeproc{ref-pedersen2024ggraph}{\textbf{26}}}

Using the same underlying output igraph object from Ball Mapper, we
created the same graph in ggraph but changed the visual representation
and added another helpful piece of information, edge weight, that was
given but not used in the default BM graph. Looking at
Figure~\ref{fig-newGraph}, we have the standard Ball Mapper graph output
on the left (Figure~\ref{fig-newGraph-1}), and our new ggraph output on
the right (Figure~\ref{fig-newGraph-2}). The most prominent changes are
the spread of the network and coloration.

Looking at Figure~\ref{fig-newGraph-2}, having the nodes spread out
helps us see the structure of our graph more clearly as well as the edge
weight between nodes. We see that highlighting the edge strength
presents a potential insight between certain nodes. Although our data is
relatively small, consisting of 65 rows, one could imagine the advantage
Figure~\ref{fig-newGraph-2} might have with a data set much larger than
this one (e.g.~looking at our time frame on a quarterly or monthly
basis). Additionally, we changed the color palette for accessibility and
readability.

For Figure~\ref{fig-newGraph-2} and Figure~\ref{fig-EDA}, we use Year as
our initial coloring variable. We started with this variable because
D\l{}otko et
al.\textsuperscript{\citeproc{ref-dlotko2019mapping}{\textbf{5}}}
provided a framework for this type of analysis. In this paper the
authors started by coloring their graphs by Year to initially understand
the shape of their data. Later, in the second half of their analysis,
they produced more graphs each colored by the point cloud's variables.
Since we have 16 dimensions to our point cloud, an analysis of that size
is outside the scope of this paper. An additional note about D\l{}otko
et al., when they colored by the variable `Year', it was not a variable
in their point cloud. Mentioned in the \emph{Introduction}, we
recognized the usefulness of this technique and extended this idea to
our point cloud. Specifically, we calculated variables from our data
that we thought represented different economic lenses: the Real Interest
Rate (Figure~\ref{fig-Fisher}), PPI-CPI Gap
(Figure~\ref{fig-PPI-CPI_Gap}), and Real Wage Growth
(Figure~\ref{fig-Wage_Growth}).

The Fisher Equation (Equation~\ref{eq-fisherEquation}) is a mathematical
relationship to approximate the Real Interest Rate. It captures the
lived experience of purchasing power for producers and consumers by
showing what is left of a nominal interest rate after accounting for
inflation--in our case we are using the nominal Federal Funds Rate (Fed
Rate). When the real interest rate is negative, inflation is higher than
the Fed Rate and is sometimes described as the Federal Reserve (the Fed)
being, ``behind the
curve.''\textsuperscript{\citeproc{ref-jacoby2022behind}{\textbf{27}}}
Being behind the curve can stem from various factors but some common
interpretations are, the Federal Reserve does not want to try and tame
inflation because of the possible negative impact on the labor market,
or the Federal Reserve is hopeful that the economy will grow into the
inflation, thereby raising wages up to or beyond inflation. On the other
hand, when the real interest rate is positive, the Fed Rate is higher
than inflation and could signal that the Federal Reserve is cautious
about inflation's impact on the economy. This is described as being
``ahead of the curve.'' If the Fed Rate is ahead of the curve, it
indicates the cost of capital is higher causing businesses to ``feel''
the rate. In reaction to business being more expensive to conduct,
employers might cut costs or freeze hiring, directly affecting the labor
market as a result of keeping prices stable. An important note about the
Fisher equation: the approximation does not hold well when nominal
interest rates (Fed Rate) are relatively high and when time spans are
short.\textsuperscript{\citeproc{ref-michaelides2024fisher}{\textbf{7}}}

\begin{equation}\phantomsection\label{eq-fisherEquation}{
\text{Real Interest Rate} \approx \text{Fed Rate} - \text{Inflation}
}\end{equation}

We calculated the PPI-CPI Gap by taking the difference between PPI -
Finished Goods and CPI (inflation). As a quick reminder, PPI represents
an ``at-cost'' revenue (before adding margin, transportation,
\emph{etc.}) that producers receive, and CPI is a common measure of
consumer inflation. It is important to note that although CPI is used
for general consumer inflation, the Federal Reserve prefers Personal
Consumption Expenditures (PCE) because it tracks more normal consumer
behavior such as substitution of products. The predominant question for
this measure was: if producers or consumers are experiencing inflation,
can our graph show the primary driving force behind the complex nature
that is inflation?

Similar in nature to the PPI-CPI gap, we took the difference between
Compensation and Inflation as an approximate measure for Real Wage
Growth. With this measure, we were interested in the relationship
between people's annual compensation and inflation, i.e., are wages
keeping up with inflation. Generally, for this relationship, positive
numbers indicate that wages are growing faster than inflation, and vice
versa if it is negative--effectively people are taking a pay cut.

\section{RESULTS}\label{results}

\emph{Exploratory Data Analysis}

Figure~\ref{fig-EDA} shows three different connected components and many
nodes not connected to any others. A connected component is when at
minimum, two nodes share one edge. If a node is not connected to any
other node, they are called singletons. When a node is a singleton, it
suggests that node is unlike all the other data points and indicates
that it is a possible outlier in the
dataset.\textsuperscript{\citeproc{ref-dlotko2019mapping}{\textbf{5}}}
Since our graph consists of 40 nodes and three components, we will
provide all the singletons in a table
(Table~\ref{tbl-Singleton_Summary}). For the three connected components,
we see there is one large and two small components. Moving clockwise,
starting with the large component on the left-hand side, we will label
them C1,\footnote{Nodes: 4, 24, 25, 27, 29, 34, 35, 36, 40}
C2,\footnote{Nodes: 7 and 8} and C3\footnote{Nodes: 2 and 28} for ease
of reference. To prime the reader, the following section briefly
discusses interpretation of Ball Mapper graphs via visual inspection. We
then apply our knowledge to our three derived metrics to finish the
section. Afterwards, we move to the \emph{Discussion} section where we
investigate the contents of our nodes and R is a necessary tool.

\begin{figure}[H]

\centering{

\includegraphics[width=0.8\linewidth,height=\textheight,keepaspectratio]{figures/fig-EDA-1.pdf}

}

\caption{\label{fig-EDA}Ball Mapper graph colored by Year; epsilon set
to 0.474.}

\end{figure}%

When doing a visual inspection of Ball Mapper graphs, there are
essentially two categories to view it through: coloration and shape.
Most often, the coloration of the Ball Mapper graph is used to highlight
a variable of interest, or certain known structures in the data. This
could take the form of highlighting known outliers, showing a coloration
pattern (e.g.~a smooth gradient from one end of the graph to the other),
or a combination of the two; i.e., nodes which do not conform to
patterning. Figure~\ref{fig-EDA} shows us no global coloration pattern,
but does reveal four sub-groups of colors in C1 as well as the
mismatching coloration for C3. From visual inspection, C1's four
sub-groups of colors show distinct time periods: 1990s (magenta),
mid-2000s (yellow), 2020s (orange), and 2000 (light terracotta). We also
notice that for C3, 2000 and 1980 are what time periods its nodes show,
while C2 shows a consistent mid-1980s coloration. Looking at shape and
structure of the graph, as mentioned above, nodes contain an edge
between them when there is an intersection of at least one point.
Additionally, the edge weight is commensurate with how many points are
in the intersection, and the overall size of the node indicates the
number of points it contains. Particular to C1, nodes which are not
directly connected by an edge (e.g.~Nodes 24 and 36), can be taken as
similar in some fashion, but also signal a significant difference within
the
component.\textsuperscript{\citeproc{ref-dlotko2019mapping}{\textbf{5}}}

Examining C2 and C3 closer, C2's nodes are colored very similarly, and
are of the same size. As a relative measure, using
Table~\ref{tbl-Singleton_Summary}, we look at our singletons and notice
that many only hold one point. Thus, comparing C2 and C3 nodes, we
should expect around one to two points in each of the nodes. C2's nodes
contain the years 1968 and 1969 in nodes 7 and 8, respectively, and 2000
is in both, creating the edge. Mentioned previously, C3's coloration is
not consistent. Node 2 contains 1961 and node 28 contains 2001 with 2002
being our bridge year. What is an interesting observation that will come
up again when looking at C1 is how different, but grouped together the
years are. This brings up questions as to why 1961 and 2001 would be
seen as similar and is explored in the \emph{Discussion} section.

Structure-wise, C1 is the largest component and contains nine total
nodes. We mentioned already that it contains three subgroups of colors
and observed the indicated years they visually represent. Of the nine
nodes that make up C1, we are provided two distinct parts: a large
triangle made of the largest nodes with strongly connected nodes (4, 29,
and 36); and a sweeping arm to the left of the triangle that contains
leaves on each end (nodes 25 and
34).\footnote{The term leaf (leaves) could also be called a flare. How we define flare is two or more nodes in a row that are hanging off the main body of a component.}
The combination of structure along with the subsets of the colors
immediately present a notable result. On the upper half of our arm we
see nodes 25, 24, and 27 all colored similarly. We have a similar
outcome for the bottom half of the arm and one node from our triangle.
Conversely, and equally interesting, the triangle's nodes are three
distinct colors, while also being the largest nodes in the graph and
containing the strongest edge strengths between them.

\begin{table}[H]

\centering{

\fontsize{8.0pt}{10.0pt}\selectfont
\begin{tabular*}{1\linewidth}{@{\extracolsep{\fill}}clclclclclcl}
\toprule
Node & Year(s) & Node & Year(s) & Node & Year(s) & Node & Year(s) & Node & Year(s) & Node & Year(s) \\ 
\midrule\addlinespace[2.5pt]
{\bfseries 1} & 1960 & {\bfseries 10} & 1971 & {\bfseries 15} & 1976, 1977 & {\bfseries 20} & 1982 & {\bfseries 30} & 2007 & {\bfseries 38} & 2021 \\ 
{\bfseries \cellcolor[HTML]{F2F2F2}{3}} & {\cellcolor[HTML]{F2F2F2}{1962, 1965}} & {\bfseries \cellcolor[HTML]{F2F2F2}{11}} & {\cellcolor[HTML]{F2F2F2}{1972, 1998}} & {\bfseries \cellcolor[HTML]{F2F2F2}{16}} & {\cellcolor[HTML]{F2F2F2}{1978}} & {\bfseries \cellcolor[HTML]{F2F2F2}{21}} & {\cellcolor[HTML]{F2F2F2}{1983}} & {\bfseries \cellcolor[HTML]{F2F2F2}{31}} & {\cellcolor[HTML]{F2F2F2}{2008}} & {\bfseries \cellcolor[HTML]{F2F2F2}{39}} & {\cellcolor[HTML]{F2F2F2}{2022}} \\ 
{\bfseries 5} & 1966 & {\bfseries 12} & 1973 & {\bfseries 17} & 1979 & {\bfseries 22} & 1984 & {\bfseries 32} & 2009 & {\bfseries } &  \\ 
{\bfseries \cellcolor[HTML]{F2F2F2}{6}} & {\cellcolor[HTML]{F2F2F2}{1967}} & {\bfseries \cellcolor[HTML]{F2F2F2}{13}} & {\cellcolor[HTML]{F2F2F2}{1974}} & {\bfseries \cellcolor[HTML]{F2F2F2}{18}} & {\cellcolor[HTML]{F2F2F2}{1980}} & {\bfseries \cellcolor[HTML]{F2F2F2}{23}} & {\cellcolor[HTML]{F2F2F2}{1985, 1986}} & {\bfseries \cellcolor[HTML]{F2F2F2}{33}} & {\cellcolor[HTML]{F2F2F2}{2010}} & {\bfseries \cellcolor[HTML]{F2F2F2}{}} & {\cellcolor[HTML]{F2F2F2}{}} \\ 
{\bfseries 9} & 1970 & {\bfseries 14} & 1975 & {\bfseries 19} & 1981 & {\bfseries 26} & 1991 & {\bfseries 37} & 2020 & {\bfseries } &  \\ 
\bottomrule
\end{tabular*}

}

\caption{\label{tbl-Singleton_Summary}Singleton nodes and the years they
contain.}

\end{table}%

Although we would need to investigate what points make up the nodes to
gain more insight into our data, on the surface what we can conclude
about C1 is: the sweeping arm seems to have two groups representing two
time periods in our range of years. We also notice that nodes 27, 35,
and 34 make their own small cycle of nodes but are a mix of the two
aforementioned time periods. Lastly, our triangle has a combination of
factors which warrant further investigation: the nodes are the largest
in size indicating they contain the most points; each pair of the nodes
(4, 29, 36) contain the strongest edge strengths of the whole graph; all
three nodes are distinct colors indicating separate time periods; and
finally, the distinct colors seem to contain a path through time going
from 2000 (node 4) around in a counter-clockwise direction to present
day (2020--2024).

\emph{Real Interest Rate}

Figure~\ref{fig-Fisher} shows us that colors representing blue/mint
green are negative values which represent inflation far outpacing the
Fed Rate. In contrast, values colored yellow/orange show us that Fed
Rates are very high compared to inflation. In general we notice that
there are not many nodes which are blue or mint colored, and most nodes
range from purple to yellow (approximately -2 to 4). Looking at the
extreme ends (less than -4 or greater than 6), we see there are about
four nodes for each end. On the low end are nodes 39, 14, 38, and 34;
and on the high end are nodes 19--21. The reader may have noticed both
groups of nodes seem to be consecutive (or close to), this is not a
random event but a feature in Ball Mapper. We touched on this in the
Methods and Procedures: Ball Mapper section, and will revisit this idea
in our \emph{Discussion}. Lastly, and the only case of a coloration
gradient, C1 shows a smooth transition from negative values at the
bottom (34) to positive in a sweeping fashion like our arm.

\emph{PPI-CPI Gap}

Figure~\ref{fig-PPI-CPI_Gap} shows a very different distribution of
colors compared to Figure~\ref{fig-Fisher}. Overall, a majority of our
nodes are below the midpoint of the legend. This shows that regardless
of the driver behind inflation on both consumers and producers, consumer
inflation tends to outpace producers. Additionally, our legend is not
balanced like it is in Figure~\ref{fig-Fisher}; this supports our
clustering of values below the middle of our legend. In fact, we see
that many nodes seem to be from -2 to -1. With a discerning eye, we
estimate to have three nodes which are at the very end of our negative
values (nodes 23, 37, and 9). However, on the opposite end there is only
one node that appears to be an extreme positive, node 39.

\begin{figure}[H]

\centering{

\includegraphics[width=0.8\linewidth,height=\textheight,keepaspectratio]{figures/fig-Fisher-1.pdf}

}

\caption{\label{fig-Fisher}Ball Mapper graph colored by the Real
Interest Rate; epsilon set to 0.474.}

\end{figure}%

\emph{Real Wage Growth}

Similar to Figure~\ref{fig-PPI-CPI_Gap}, there is a clear concentration
of values on one end of our legend Real Wage Growth shown in
Figure~\ref{fig-Wage_Growth}. Using the center of the legend as a
reference point, we can see that Figure~\ref{fig-Wage_Growth} is even
more skewed than Figure~\ref{fig-PPI-CPI_Gap}; its zero starts two tick
marks from the center instead of one. To give a range for the
overwhelming yellow coloring, most values seem to fall between -2.5 and
2.5, with most of these values lying right around zero. With many values
indicating zero, our selected time period tells us that wages have
seemed to keep pace with inflation. Although
Figure~\ref{fig-Wage_Growth} and Figure~\ref{fig-PPI-CPI_Gap} are
showing different measures, Real Wage Growth shows a much longer tail.
Intuitively this makes sense because with the PPI-CPI Gap, both
producers and consumers play a similar role, i.e., they are a consumer
to someone. On the other hand, if inflation spikes and a business cannot
keep up with wages, there can be much more of a divergence. That said,
nodes 18, 19, and 17 might be cases for this and further investigation
is found in the \emph{Discussion} section.

\section{DISCUSSION}\label{discussion}

Often with Ball Mapper analyses EDA and visual inspection are the first
steps, and then inspection of what points are housed in each node done
after. This process of looking into each node helps us understand why we
are seeing certain structures and the coloration. To do this
investigative work, we combined all of our dimensions in our point cloud
(Table~\ref{tbl-Analysis_Data}) with our three derived economic measures
(Real Interest Rate, PPI-CPI Gap, Real Wage Growth), as well as the
Year, GDP, NBER recession/expansion indication, and a Row ID to join all
the data together (Table~\ref{tbl-Discussion_Data}). For the remainder
of the discussion we will look closely at each component, starting with
C2 and C3. We will spend a significant portion on C1 due to its size and
complexity. Then we will turn to the singletons and finish with some
questions and future research.

\emph{C2}

C2 is similarly colored across nodes, but this is misleading because of
the average value used for coloration. Nodes 7 and 8 contain the years
1968 and 1969, respectively, and the year 2000 forms the connecting
edge. Ball Mapper has paired these years together because of their
shared macroeconomic structures. What our analysis suggests is that C2
represents an economic peak that shows a demand-pull inflation
structure, reflecting the adage ``too much money chasing too few
goods.''\textsuperscript{\citeproc{ref-investopedia_demandpull}{\textbf{28}}}
Further, C2 describes a clean or structurally coherent peak, and shows
that this kind of peak contains distinct stages.

In Figure~\ref{fig-Fisher} through Figure~\ref{fig-Wage_Growth}, nodes 7
and 8 both show real interest rates approaching 3\%, PPI-CPI gaps close
to zero, as well as near-zero values for real wage growth. These three
metrics indicate the economy was running hot but stable. Prices were not
rising disproportionately for either producers or consumers. Both nodes
also show strong GDP growth and very tight labor markets where
unemployment has been falling through the points in C2. Where C2 splits
into stages is that node 7 seems to represent early-stage demand-pull
inflation, and node 8 represents late-stage. In node 7 this shows up in
the positive housing starts (5.6\%), real wage growth (1.3\%), moderate
inflation (3.7\%), and a Fed Rate (6.0\%) that has been rising
throughout the years. Node 8, however, tells a different story: housing
starts roll over (-2.9\%), real wage growth shrinking to almost zero
(0.17\%), inflation continuing to rise (4.3\%) and the Federal Reserve
aggressively responding with a nominal rate of 7.2\%.

The parallel between early and late-stage demand-pull inflation is
reflected in distinct historical policy and business choices. The
late-1960s demand was primarily supported by mid-1960s fiscally driven
policies, spending by the government on the Vietnam war, expansion of
the workforce from college-aged graduates, and family
formation.\textsuperscript{\citeproc{ref-marsh2021inflation}{\textbf{29}}}
In 2000, corporate spending accelerated in anticipation of the digital
age, playing a similar demand-driving role to the 1960s policies. This
inevitably fueled the ``wealth effect'', encouraging consumers to drive
up
demand.\textsuperscript{\citeproc{ref-evercore2024booms}{\textbf{30}},
\citeproc{ref-sussman2019wealth}{\textbf{31}}}

\begin{figure}[H]

\centering{

\includegraphics[width=0.8\linewidth,height=\textheight,keepaspectratio]{figures/fig-PPI-CPI_Gap-1.pdf}

}

\caption{\label{fig-PPI-CPI_Gap}Ball Mapper graph colored by PPI less
CPI (PPI-CPI Gap); epsilon set to 0.474.}

\end{figure}%

\emph{C3}

Similar to C2, node 2 is misleading on the surface in
Figure~\ref{fig-EDA}, it contains years 1961 and 2002 but the average
value places it around the 1980s. Node 28 however is much closer in
years, 2001 and 2002. Between the two nodes, though, we notice there is
a common theme: they show an economy which is at a trough and has
stabilized. What we mean is that this component seems to represent when
a contraction has ended, but recovery has not started. Further, unlike
C2 we find that C3 shows two different types of troughs: labor market
trough and an investment trough. Both nodes show average muted GDP
growth (2.2\% for node 2; 1.4\% for node 28), low Fed Rates (both nodes
less than 2.8\%) falling throughout the year, housing starts are
positive, and inflation is under control (1.3\% and 2.2\%). However,
where they differ is that the labor market seems to break in node 2.
Here we see personal income and compensation relatively intact (1.7\%
and 1.0\% respectively), and high unemployment (6.2\%) that has been
rising being cushioned by elevated transfer utilization. Conversely,
node 28 also has elevated unemployment levels, but personal income and
compensation are low, causing real wage growth to turn negative
(-2.0\%). Additionally, domestic investment pulls back aggressively
(-3.3\%), suggesting business pulling back is the primary driver behind
this node.

\emph{C1}

As mentioned above, what is striking about C1 is that there are two
shapes: a triangle and a sweeping arm that runs along the triangle's
left side. Although we have not found any literature that explores the
relationship between nodes and edge strength (how many points they
share), the triangle (denoted body from now on) was something
immediately of interest. Beginning with Real Wage Growth, it falls
between -2.5\% and 2.5\% (Figure~\ref{fig-Wage_Growth}), with an average
value of -0.065. A value close to zero tells us that for the points in
C1, wages have generally kept up with inflation. The PPI-CPI Gap
(Figure~\ref{fig-PPI-CPI_Gap}) shows there is a main group of nodes
centered around negative one with three nodes higher (25, 34, and 29),
and one lower (40). These mostly negative values suggest that consumer
inflation is generally outpacing producers in C1. It's worth noting,
however, that our PPI measure does not consider profit margin added,
transportation costs included, or other on-the-shelf costs. So this
small differential that we see in C1 possibly could be reflecting that
gap in pricing, but a more in depth analysis or additional data in our
point cloud is needed to provide a more definitive answer. The real
interest rate graph, Figure~\ref{fig-Fisher}, displays a color gradient
for C1 that immediately stood out to us. Seeing that it moves from
negative to positive values in a semi-orderly manner supports the idea
that the nodes with negative values show inflation was outpacing the
nominal Fed Rate; and as we move toward the positive values, we see that
relationship flip, with the Fed Rate exceeding inflation. Supposing that
the Federal Reserve was behind the curve for the negative values and
ahead for positive, along with our knowledge that overall real wages
have kept up and the PPI-CPI Gap is relatively small, we start to get a
picture of what C1 could represent.

\begin{figure}[H]

\centering{

\includegraphics[width=0.8\linewidth,height=\textheight,keepaspectratio]{figures/fig-Wage_Growth-1.pdf}

}

\caption{\label{fig-Wage_Growth}Ball Mapper graph colored by
Compensation less Inflation (Real Wage Growth); epsilon set to 0.474.}

\end{figure}%

Indeed, what we have found for C1 is that it represents the U.S. economy
at different states of growth. This does not mean that every node in C1
is an instance of good economic conditions, but rather nodes represent
certain ``phases'' that the economy goes through until it turns into a
state of overheating expansion (peak) or coming out of a state of
contraction (trough). What we have also found is that these peaks and
troughs are not included in C1, but are represented by C2 and C3
respectively.

\emph{C1: Recovery and Resistance}

Broadly, the stages can be categorized in different phases: recovery (34
and 35), stability (29, 4, and 36), resistance (24, 25), and
non-canonical (27 and 40). Starting with our recovery phase, we find
signs of an economy that has been beaten down but is starting to come
back to life. We find that 34 shows a very fragile economy and 35, being
connected to more nodes, indicates more stability. The first thing we
notice is that housing starts is positive, so builders are feeling
confident that they can start projects. The ``big four'' economic
indicators--GDP, Inflation, Unemployment, Fed Rate--on average, show
signs of improvement. GDP is increasing, inflation cooling, Unemployment
starting to fall, and the Fed Rate relatively low because they are
trying to stimulate the economy. Additionally, real interest rates are
negative and real wage growth is improving, suggesting that producers
and consumers feel more confident about the economy.

Next, we have what we call the resistance phase. Unlike our recovery
phase, nodes 24 and 25 show two different sides of resistance. Node 24
shows an economy which feels like it is still growing and running well.
Unemployment, on average, is 5.4\% and trending downwards, inflation
seems to be under control, around 3.1\%. However, at this stage there
seems to be some cracks forming in the economy, the first being housing
starts turning negative. What is also worrisome is the high Fed Rate
(5.9\%) and the negative PPI-CPI gap (-1.1\%). The combination of these
suggests that the Federal Reserve sees that CPI is running ahead of PPI
and are concerned about producer inflation catching up to consumer
inflation. With CPI outpacing PPI, consumer prices are currently rising,
however if producer costs start to rise, the PPI-CPI Gap will start to
close and inflation then is a problem for both producers and consumers.
To counteract this, the Federal Reserve is preemptively trying to kill
inflation for consumers with higher rates. Unfortunately, this has
knock-on effects to the labor market restricting wage growth and
possibly reversing the direction of unemployment. Node 25 shows the
other side of what 24 was warning us about. Housing slows down even more
(from close to -5\% to -10\%, on average), and the PPI-CPI Gap shrinks
to almost zero; unlike node 24, producer price inflation is now catching
up to consumer price inflation. Since both producers and consumer prices
are rising, the Federal Reserve continues to raise rates to an average
7\%, sending average real interest rates above 3\% and risking the labor
market to decline and the economy to possibly enter a recession. We also
see real wage growth plummeting from -0.16\% (node 24) to -2.23\% (node
25). To illustrate this differential, suppose we were given \$1,000 at
the end of the year and to either had to give some back (negative wage
growth), or we received extra (positive wage growth). In node 24 we get
that \$1,000 and give back \$1.60, however in 25 we would have to give
back \$22.30. So, node 25 shows a ``last resort'' node before turning
into a singleton node. This is similar to 34 in the sense that it is
only connected to the other nodes in C1 by one other node, supporting
this idea that they are the exit and entrance into economic growth.

\begin{table}[H]

\centering{

\fontsize{8.0pt}{10.0pt}\selectfont
\begin{tabular*}{1\linewidth}{@{\extracolsep{\fill}}llll}
\toprule
 &  &  &  \\ 
\midrule\addlinespace[2.5pt]
Year & PPI Finished Change & Personal Income & Compensation \\ 
{\cellcolor[HTML]{F2F2F2}{Entrepreneurship}} & {\cellcolor[HTML]{F2F2F2}{Transfers}} & {\cellcolor[HTML]{F2F2F2}{GDP}} & {\cellcolor[HTML]{F2F2F2}{Consumption}} \\ 
Domestic Investment & Government Spending & Exports & Imports \\ 
{\cellcolor[HTML]{F2F2F2}{Inflation}} & {\cellcolor[HTML]{F2F2F2}{Unemployment}} & {\cellcolor[HTML]{F2F2F2}{Unemployment Change}} & {\cellcolor[HTML]{F2F2F2}{Fed Rate}} \\ 
Fed Rate Change & Housing Change & Recession Span (0-1 by 0.25) & Expansion Span (0-1 by 0.25) \\ 
{\cellcolor[HTML]{F2F2F2}{Row Number}} & {\cellcolor[HTML]{F2F2F2}{Real Interest Rate}} & {\cellcolor[HTML]{F2F2F2}{PPI CPI Gap}} & {\cellcolor[HTML]{F2F2F2}{Real Wage Growth}} \\ 
\bottomrule
\end{tabular*}

}

\caption{\label{tbl-Discussion_Data}Data used to review nodes, included
additional data that was excluded from our point cloud.}

\end{table}%

\emph{C1: Stability}

The final phase tells a story of stability and contains three stages or
variations. Starting with node 29, we label it emerging-stability; it
shows an economy that has found its legs and is now starting to gain
some momentum. Across the three stages, 29 has the highest average
unemployment at 5.4\% and the lowest average Fed Rate (1.5\%). The
Federal Reserve (the Fed) is intentionally keeping rates low, making the
real interest rate -0.81\%, to stimulate the economy. The goal is to
make borrowing cheaper to create an incentive to producers and consumers
to spend. The PPI-CPI gap at 0.16\% informs us that producer prices are
slightly growing faster than consumer prices. There are no clear reasons
as to why prices are higher on the producer side but some possibilities
in the context of emerging-stability are: producers are seeing a pick up
in demand or retailers that sell to consumers have not passed on costs
to consumers yet in anticipation that conditions get better. Lastly,
real wage growth is negative (-0.27\%) and consumption is at a moderate
3.08\%--the middle average of the three nodes. These two averages
support the signs that consumer confidence is gaining and the economy is
coming back to life.

The story of node 36 is one of extended-stability. Between the three
stages, the producers and consumers seem to be more cautious; on
average, the labor market is tightening (unemployment at 4.2\%), and GDP
(2.5\%) along with housing starts (2.4\%) are at their lowest. We also
see consumer spending slowing down between our three nodes at 2.56\%,
which is consistent with real wage growth being close to zero (0.24\%).
The Federal Reserve at this stage has achieved its goal shown by real
interest rates at 0.01\% and inflation at 2.10\%. With its goals met
(stable prices and maximum employment), it shifts to actively monitoring
for any signs expansion is starting to overheat.

Unlike node 36, node 4 shows the Federal Reserve comfortable with how
the economy is running and is letting the economy run its course; we
have identified it as full-stability. What gives the Federal Reserve
confidence for a ``hands off'' approach is the low inflation rate
(1.94\%), unemployment at 5.03\%, moderate housing starts (3.96\%), and
strong GDP growth of 3.72\%. Looking at our three derived metrics, our
real interest rate is 0.86\%, indicating a slight drag on borrowing,
allowing the economy to moderate itself as it continues to run
efficiently. Our PPI-CPI gap (-0.82\%) is telling us consumer prices are
rising faster than producer prices. This gap in combination with the
fact that this stage has the highest consumption (3.76\%) informs
producers that demand is strong at the retail level. Consumers also
benefit from real wage growth at 1.56\%; their wages are outpacing
inflation, giving them an increase in purchasing power. To illustrate
using our example for node 25, if we got \$1,000 at the end of the year,
this positive real wage growth says that we get an extra \$15.60.

\emph{C1: Non-canonical}

Up to this point in our discussion we have seen various stages of the
U.S. economy's expansion behavior. Two of our nodes in C1, 40 and 27, do
not align with any ``phase'' we have talked about thus far. Starting
with node 40, it is connected to both node 4 and 36; to recall, node 4
is our full-stability and 36 is extended-stability. What we find with
node 40 is that consumption is the driving factor. Despite headwinds to
the consumer--inflation at 3.48\% and real wage growth at -1.38\%--they
continue to buy. However, housing starts have turned negative (-6.20\%),
signaling that builders are stepping back and consumers are stepping up.
The Federal Reserve, observing consumer resilience and builders'
caution, is trying to cool the economy (real interest rates at 1.61\%)
without crashing the tight labor market (unemployment at 3.83\%).

Why we chose to categorize node 27 as non-canonical is because it is the
only node which is connected to all three of our phases. Looking at
average values, node 27 shows a balancing act of trying to stimulate an
economy that is slow to recover without further entrenching sticky,
elevated unemployment. On the surface there are signs of confidence. GDP
is at 3.33\%, there is high domestic investment (8.00\%), and housing
starts have picked up (10.78\%). When looking at the real interest rate,
the Fed Rate, and unemployment, though, we see a Federal Reserve trying
not to spur inflation while trying to bring unemployment down. On one
hand we have a real interest rate of 1.01\% that puts a small drag on
growth, but on the other the Fed Rate has changed by -1.07\% and
unemployment is elevated (6.60\%) but unchanging (-0.04\%).

\emph{Singletons}

So far we have discussed how our three components represent the various
growth stages, structural peaks, and structural troughs for the U.S.
economy. As one might expect, singleton nodes do represent states of
contraction. However, in addition to the contractions, we found that
singletons also represent non-structural, or unstable growth that could
feel contractionary. To show that our singletons contain economic
contractions we will use a combination of known economic crises as well
as observations we found from our derived economic measures. We will
also weave in time periods where the economy was growing but in an
unusual or unstable way.

Perhaps the most obvious years to show are The Great Recession
(2007--2009). Although The Great Recession is known to be '07--'09,
nodes 30--32 are the years 2007--2009; however, node 33, another
singleton, is the year 2010 seen in Table~\ref{tbl-Singleton_Summary}.
These years are marked by deep, negative housing starts, very large Fed
Rate cuts, and high unemployment. Another more recent time period was
the COVID-19 Pandemic. Starting in 2020, the pandemic lasted for a few
years, and we see that nodes 37--39 harbor these years: 2020, 2021, and
2022. In 2020, we had a contraction with unemployment rising very
quickly and substantial transfers from government stimulus. Following
2020, we had economic expansion in 2021 and 2022 despite people feeling
the economy was in a state of contraction. In these two years, we see a
pattern of growth that feels bad: positive GDP, falling unemployment,
and positive domestic investment mixed with rising inflation, shrinking
real wage growth, slowing housing starts, and rapidly rising PPI.

Looking at Table~\ref{tbl-Singleton_Summary}, something the reader might
notice is that from 1970--1986, nearly every year is a singleton with
the exception of 1976--77 and 1985--86. This tells us that Ball Mapper
has found meaningful differences between these years such that no year
is similar in economic structure to another. Throughout this 16-year
period, a few factors that contributed to the distinctiveness are
external supply shocks and monetary policy decisions. The oil embargo of
1973--74 was a textbook case of cost-push inflation. This supply shock
is visible in Figure~\ref{fig-Wage_Growth}, where real wage growth
reaches one of its lowest values. During 1974 workers' real wages fell
to -12.06\%. Returning to our previous example, this is receiving that
\$1,000 but having to give back \$120.60 of
it.\textsuperscript{\citeproc{ref-corbett2013oilshock}{\textbf{32}}}

In 1979 Paul Volcker became the Federal Reserve chairman. Prior to
Volcker, the U.S. economy was dealing with the stagflation from the
1970s, characterized by high inflation, rising unemployment, and
stagnant growth. So, to eliminate the climbing inflation, Volcker
decided to aggressively hike rates with rates peaking in 1981.
Specifically, at their peak, Volcker raised the nominal interest rate
(Fed Rate) to 16.39\% at the expense of the labor market. We see this
decision distinctly in Figure~\ref{fig-Fisher}. Nodes 19--22 represent
the years 1981--1984, which show us another example of a contractionary
and expansionary period, this time through monetary policy. In 1981,
real interest rates were 6.54\%, showing that the Federal Reserve was
ahead of the curve. By doing this, Volcker achieved his goal of lowering
inflation from 9.87\% in 1981 to 5.98\% in 1982. Raising the Fed Rate,
though, made the cost of capital very expensive which caused
unemployment to rise and stay around 9.65\% from 1982--83. So, during
1981--82 we had a contraction period, and in 1983--84, we saw economic
growth with instability in employment and elevated Fed
Rates.\textsuperscript{\citeproc{ref-bryan2013greatinflation}{\textbf{33}}}

\newpage

\emph{Final Notes and Future Research}

As the last part of our discussion, there are two primary areas. The
first half touches on the limitations we found both with Ball Mapper and
the design of our analysis. The second examines some of the unanswered
questions that arose during the analysis which are possible avenues for
future research.

The primary limitations we had were the choice of time period and number
of variables. Initially we used a set of variables we thought would
represent the U.S. economy as a whole, but found that they did not all
have consistent time ranges. Some data series, such as the Federal Funds
Rate, ranged as far back as the early 1900s, and others were only
created recently, such as Foreign Direct Investment (2014). While we had
65 years of data, we were not able to capture early structural economic
changes or events such as the Great Depression, the New Deal, and WWII.
Alternatively, we could have chosen a different frequency of data, such
as monthly or quarterly. This would have given us more points in our
point cloud and perhaps a more nuanced picture of the U.S. economy. Our
analysis also could have included more dimensions, testing Ball Mapper's
ability to handle a multitude of dimensions. It also could have asked a
more focused economic question to see if there is a nuanced structure in
a particular area of the economy.

Focusing on the technical limitations, the first is that the standard
Euclidean Distance between points is the only distance metric used.
Currently, in Ball Mapper (R version) there is no way to easily change
the distance metric. To do so one would need to copy the source code of
the \emph{BallMapper} function and either: adapt the code to add an
argument for a distance metric, or hard code it into a copy of the
original function. Additionally, the \emph{BallMapper} function in R
uses a greedy, sequential algorithm to select ball centers. When a
landmark is chosen, the Ball Mapper algorithm does so in a sequential
top-down order. This rigid structure is good when others would like to
verify the output of the graph; having the same dataset, in the same
order, using the same epsilon, will produce the exact same graph.
However, Ball Mapper could be susceptible to seasonal data where you
have large naturally occurring spaces between the data. While
seasonality is a valid characteristic of a dataset, it might not be the
true underlying structure. Because of Ball Mapper's greedy, sequential
algorithm, it may pick up on the seasonality of the data right away
rather than comprehensively looking at the dataset. To fix some of these
issues, one could randomize the dataset every time Ball Mapper is run,
or make sure to seasonally adjust or normalize any data. Currently, some
of these fixes have been implemented in the package \emph{pyBallMapper},
written in Python. In this package they have added an order argument to
allow randomization of rows. They also have included a distance metric
argument so a user can choose what distance metric they would like to
use. Finally, \emph{pyBallMapper} uses Python's NumPy package which
utilizes vectorization, speeding up the sequential nature of the Ball
Mapper algorithm.

To start the second half, we noticed during the analysis there is not a
clear way to know whether Ball Mapper is giving us a stable graph of our
data. We discussed that if given the same data, in the same order, it
indeed will give us the same graph. However, we did not find any
literature answering this exact question. The closest paper we could
find was one which studied the Brexit Vote using Ball Mapper. At the end
of this paper they wrote about a robustness test they implemented by
showing why the vote to leave the European Union
passed.\textsuperscript{\citeproc{ref-rudkin2024brexit}{\textbf{34}}}
What is unclear is how they ran a stated 10,000 repetitions. From what
we know about the Ball Mapper algorithm, we could infer that they
randomize the order of their data every time, or took infinitesimal
slices of an interval such that theoretically Ball Mapper would produce
a different graph. The former would be a good standard to show general
robustness for any Ball Mapper output and could look like running many
repetitions while randomizing the order of the data every time, and
simultaneously collecting summary statistics of each graph. Some example
statistics might include the number of components, nodes, and edges at a
singular epsilon or a range of epsilons. One could also look at edges
and edge weights for consistency, as well as some sort of base value,
noting if there is a large or small variance.

Moving to more abstract questions that came up during our analysis, our
questions revolved around nodes' relationships to each other based on
various factors: edge strength, connectedness (degree), and node size.
Acknowledging that node relationships could have a combination of these
factors, our primary questions are: does edge strength between nodes
indicate a stronger relationship? For our graph, does the edge strength
between nodes 4 and 29 represent a stronger structural signal than the
edge between node 25 and 24? There is also the question of connectedness
or degree--how many edges a node has. Similar to edge strength, does the
fact that node 27 has three edges indicate its importance over node 34
with just one edge? Lastly, and in combination with the previous two,
does node size carry any relational value? Using our node 4 and node 29
example again, since they are the largest nodes with the strongest edge
weight, does this combination indicate that their relationship is
stronger? That is, if we were to look at this graph over many epsilons,
would the connection between nodes 4 and 29 persist while the connection
between nodes 4 and 40 does not? All of these questions came up during
our analysis, and to our knowledge have not been studied; they are all
areas for future research on Ball Mapper.

\section{CONCLUSIONS}\label{conclusions}

Although there is much further work to do on developing Ball Mapper,
through our analysis we have shown how TDA, and Ball Mapper in
particular, is able to give insight into data that we might not see or
describe easily. Throughout this paper we have hopefully guided our
reader to better understand how Ball Mapper works, as well as shown a
path to start implementing it. From our structural map of the U.S.
economy from 1960 to 2024, Ball Mapper revealed structural expansion,
structural ``clean'' peaks and troughs, as well as showing contractions
and non-organic growth. By displaying the Ball Mapper output with a new
graph, we were able to clearly explain how to systematically take
high-level, global interpretations of our Ball Mapper output and use
them in combination with lower-level, local, node-specific
investigation. We also explained some of the limitations we discovered
and touched on future work based on questions that came up for us during
the analysis.

Ball Mapper has been applied to many fields ranging across the social
sciences, finance, and the life sciences. Ball Mapper, and TDA in
general, may have an advantage in some of the social sciences because it
can combine both qualitative and quantitative data in the same analysis.
For instance, in our Brexit example from our \emph{Discussion}, the
authors sought to understand qualitative and quantitative motivations
behind the Brexit vote outcome. Some of the variables they used were
Household composition (single, married, etc.), Car ownership (number of
cars), Housing Tenure (Owned, Private rental, etc.), and Age. Whether
Ball Mapper brings new insights, or shows what is known in a unique way,
we hope that this paper serves as an accessible entry point for
researchers of all disciplines who are interested in what TDA can offer.

\section{REFERENCES}\label{references}

\begin{CSLReferences}{1}{0}
\bibitem[\citeproctext]{ref-madukpe2025mapper}
\CSLLeftMargin{\textbf{1}. }%
\CSLRightInline{Madukpe, V. N., Ugoala, B. C. and Zulkepli, N. F. S.
(2025) A comprehensive review of the {M}apper algorithm, a topological
data analysis technique, and its applications across various fields
(2007--2025), \emph{arXiv preprint arXiv:2504.09042}
\emph{\url{https://arxiv.org/abs/2504.09042}}}

\bibitem[\citeproctext]{ref-chazal2021introduction}
\CSLLeftMargin{\textbf{2}. }%
\CSLRightInline{Chazal, F. and Michel, B. (2021) An introduction to
topological data analysis: Fundamental and practical aspects for data
scientists, \emph{Front. Artif. Intell.} 4, 667963.
https://doi.org/\href{https://doi.org/10.3389/frai.2021.667963}{10.3389/frai.2021.667963}}

\bibitem[\citeproctext]{ref-singh2007topological}
\CSLLeftMargin{\textbf{3}. }%
\CSLRightInline{Singh, G., Mémoli, F. and Carlsson, G. (2007)
\href{https://www.researchgate.net/profile/Facundo-Memoli/publication/221571174_Topological_Methods_for_the_Analysis_of_High_Dimensional_Data_Sets_and_3D_Object_Recognition/links/02bfe50fe76b396b8f000000/Topological-Methods-for-the-Analysis-of-High-Dimensional-Data-Sets-and-3D-Object-Recognition.pdf}{Topological
methods for the analysis of high dimensional data sets and {3D} object
recognition}, \emph{in Eurographics symposium on point-based graphics},
}

\bibitem[\citeproctext]{ref-dlotko2019ball}
\CSLLeftMargin{\textbf{4}. }%
\CSLRightInline{Dłotko, P. (2019) Ball mapper: A shape summary for
topological data analysis, \emph{arXiv preprint arXiv:1901.07410}
\emph{\url{https://arxiv.org/abs/1901.07410}}}

\bibitem[\citeproctext]{ref-dlotko2019mapping}
\CSLLeftMargin{\textbf{5}. }%
\CSLRightInline{Dłotko, P., Rudkin, S. and Qiu, W. (2019) Topologically
mapping the macroeconomy, \emph{arXiv preprint arXiv:1911.10476}
\emph{\url{https://arxiv.org/abs/1911.10476}}}

\bibitem[\citeproctext]{ref-mcclung2024fisher}
\CSLLeftMargin{\textbf{6}. }%
\CSLRightInline{McClung, S. (2024)
\href{https://inomics.com/terms/fisher-effect-1545486}{Fisher effect}, }

\bibitem[\citeproctext]{ref-michaelides2024fisher}
\CSLLeftMargin{\textbf{7}. }%
\CSLRightInline{Michaelides, P. G. (2024)
\href{https://doi.org/10.1007/978-3-031-76140-9_9}{The fisher equation},
\emph{in 21 equations that shaped the world economy: Understanding the
theory behind the equations} 113--120, Springer, }

\bibitem[\citeproctext]{ref-bls_ppi_cpi_comparison}
\CSLLeftMargin{\textbf{8}. }%
\CSLRightInline{Bureau of Labor Statistics (2023)
\href{https://www.bls.gov/ppi/methodology-reports/comparing-the-producer-price-index-for-personal-consumption-with-the-us-all-items-cpi-for-all-urban-consumers.htm}{How
does the producer price index differ from the consumer price index?
Comparing the personal consumption {PPI} with the {CPI}}, }

\bibitem[\citeproctext]{ref-sullivan1997trends}
\CSLLeftMargin{\textbf{9}. }%
\CSLRightInline{Sullivan, D. G. (1997) Trends in real wage growth,
\emph{Chicago Fed Letter}
\emph{\url{https://www.chicagofed.org/publications/chicago-fed-letter/1997/march-115}}}

\bibitem[\citeproctext]{ref-bea_data}
\CSLLeftMargin{\textbf{10}. }%
\CSLRightInline{Bureau of Economic Analysis (n.d.)
\href{https://www.bea.gov/itable}{Interactive data application}, }

\bibitem[\citeproctext]{ref-fred_data}
\CSLLeftMargin{\textbf{11}. }%
\CSLRightInline{Federal Reserve Bank of St. Louis (n.d.)
\href{https://fred.stlouisfed.org/}{Federal reserve economic data
(FRED)}, }

\bibitem[\citeproctext]{ref-bea_r_package}
\CSLLeftMargin{\textbf{12}. }%
\CSLRightInline{Batch, A., Chen, J. and Kampas, W. (2018)
\emph{\href{https://CRAN.R-project.org/package=bea.R}{Bea.r: Bureau of
economic analysis API}}, }

\bibitem[\citeproctext]{ref-fredr_package}
\CSLLeftMargin{\textbf{13}. }%
\CSLRightInline{Boysel, S. and Vaughan, D. (2021)
\emph{\href{https://CRAN.R-project.org/package=fredr}{Fredr: An r client
for the 'FRED' API}}, }

\bibitem[\citeproctext]{ref-hughes2019value}
\CSLLeftMargin{\textbf{14}. }%
\CSLRightInline{Hughes-Cromwick, E. and Coronado, J. (2019) The value of
{US} government data to {US} business decisions, \emph{J. Econ.
Perspect.} 33, 131--146.
https://doi.org/\href{https://doi.org/10.1257/jep.33.1.131}{10.1257/jep.33.1.131}}

\bibitem[\citeproctext]{ref-dallasfed_nominal}
\CSLLeftMargin{\textbf{15}. }%
\CSLRightInline{Federal Reserve Bank of Dallas (n.d.)
\href{https://www.dallasfed.org/research/basics/nominal}{Deflating
nominal values to real values}, }

\bibitem[\citeproctext]{ref-hamilton2014logarithms}
\CSLLeftMargin{\textbf{16}. }%
\CSLRightInline{Hamilton, J. D. and Chinn, M. (2014)
\href{https://econbrowser.com/archives/2014/02/use-of-logarithms-in-economics}{Use
of logarithms in economics}, }

\bibitem[\citeproctext]{ref-census_soc_methodology}
\CSLLeftMargin{\textbf{17}. }%
\CSLRightInline{U.S. Census Bureau (n.d.)
\href{https://www.census.gov/construction/soc/methodology.html}{Survey
of construction (SOC): methodology}, }

\bibitem[\citeproctext]{ref-bls_ppi_handbook}
\CSLLeftMargin{\textbf{18}. }%
\CSLRightInline{Bureau of Labor Statistics (n.d.)
\href{https://www.bls.gov/opub/hom/ppi/home.htm}{Producer price indexes:
Handbook of methods}, }

\bibitem[\citeproctext]{ref-bls_cpi_concepts}
\CSLLeftMargin{\textbf{19}. }%
\CSLRightInline{Bureau of Labor Statistics (2025)
\href{https://www.bls.gov/opub/hom/cpi/concepts.htm}{Consumer price
index: concepts}, }

\bibitem[\citeproctext]{ref-rudkin2025ballmapper}
\CSLLeftMargin{\textbf{20}. }%
\CSLRightInline{Rudkin, S. (2025) An introduction to topological data
analysis ball mapper in r, \emph{arXiv preprint arXiv:2504.14081}
\emph{\url{https://arxiv.org/abs/2504.14081}}}

\bibitem[\citeproctext]{ref-qiu2020refining}
\CSLLeftMargin{\textbf{21}. }%
\CSLRightInline{Qiu, W., Rudkin, S. and Dłotko, P. (2020) Refining
understanding of corporate failure through a topological data analysis
mapping of {A}ltman's {Z}-score model, \emph{Expert Syst. Appl.} 156,
113475.
https://doi.org/\href{https://doi.org/10.1016/j.eswa.2020.113475}{10.1016/j.eswa.2020.113475}}

\bibitem[\citeproctext]{ref-ballmapper_package}
\CSLLeftMargin{\textbf{22}. }%
\CSLRightInline{Dłotko, P. (2019)
\emph{\href{https://CRAN.R-project.org/package=BallMapper}{BallMapper:
The ball mapper algorithm}}, }

\bibitem[\citeproctext]{ref-bls_cps_concepts}
\CSLLeftMargin{\textbf{23}. }%
\CSLRightInline{Bureau of Labor Statistics (2018)
\href{https://www.bls.gov/opub/hom/cps/concepts.htm}{Current population
survey: concepts}, }

\bibitem[\citeproctext]{ref-fed_monetary_goals}
\CSLLeftMargin{\textbf{24}. }%
\CSLRightInline{Board of Governors of the Federal Reserve System (2021)
\href{https://www.federalreserve.gov/monetarypolicy/monetary-policy-what-are-its-goals-how-does-it-work.htm}{Monetary
policy: What are its goals? How does it work?}, }

\bibitem[\citeproctext]{ref-csardi2006igraph}
\CSLLeftMargin{\textbf{25}. }%
\CSLRightInline{Csardi, G. and Nepusz, T. (2006) The igraph software
package for complex network research, \emph{InterJournal, Complex
Systems} 1695, 1--9. \emph{\url{https://igraph.org}}}

\bibitem[\citeproctext]{ref-pedersen2024ggraph}
\CSLLeftMargin{\textbf{26}. }%
\CSLRightInline{Pedersen, T. L. (2024)
\emph{\href{https://CRAN.R-project.org/package=ggraph}{Ggraph: An
implementation of grammar of graphics for graphs and networks}}, }

\bibitem[\citeproctext]{ref-jacoby2022behind}
\CSLLeftMargin{\textbf{27}. }%
\CSLRightInline{Jacoby, T. (2022)
\href{https://www.nasdaq.com/articles/what-does-it-mean-when-the-fed-is-behind-the-curve}{What
does it mean when the fed is 'behind the curve?'}, }

\bibitem[\citeproctext]{ref-investopedia_demandpull}
\CSLLeftMargin{\textbf{28}. }%
\CSLRightInline{Investopedia (n.d.)
\href{https://www.investopedia.com/terms/d/demandpullinflation.asp}{Demand-pull
inflation: Definition, how it works, causes, vs. Cost-push inflation}, }

\bibitem[\citeproctext]{ref-marsh2021inflation}
\CSLLeftMargin{\textbf{29}. }%
\CSLRightInline{Marsh, D. (2021) Look at 1960s, not 1970s, to learn how
{US} inflation took hold, \emph{OMFIF}
\emph{\url{https://www.omfif.org/2021/11/look-at-1960s-not-1970s-to-learn-how-us-inflation-took-hold/}}}

\bibitem[\citeproctext]{ref-evercore2024booms}
\CSLLeftMargin{\textbf{30}. }%
\CSLRightInline{Evercore Wealth and Trust (2024)
\href{https://evercorewealthandtrust.com/booms-and-busts-a-brief-history-of-capital-spending-cycles/}{Booms
and busts: A brief history of capital spending cycles}, }

\bibitem[\citeproctext]{ref-sussman2019wealth}
\CSLLeftMargin{\textbf{31}. }%
\CSLRightInline{Sussman, A. L. (2019) New estimates of the stock market
wealth effect, \emph{NBER Digest}
\emph{\url{https://www.nber.org/digest/aug19/new-estimates-stock-market-wealth-effect}}}

\bibitem[\citeproctext]{ref-corbett2013oilshock}
\CSLLeftMargin{\textbf{32}. }%
\CSLRightInline{Corbett, M. (2013)
\href{https://www.federalreservehistory.org/essays/oil-shock-of-1973-74}{Oil
shock of 1973--74}, }

\bibitem[\citeproctext]{ref-bryan2013greatinflation}
\CSLLeftMargin{\textbf{33}. }%
\CSLRightInline{Bryan, M. (2013)
\href{https://www.federalreservehistory.org/essays/great-inflation}{The
great inflation}, }

\bibitem[\citeproctext]{ref-rudkin2024brexit}
\CSLLeftMargin{\textbf{34}. }%
\CSLRightInline{Rudkin, S., Barros, L., Dłotko, P. and Qiu, W. (2024) An
economic topology of the {B}rexit vote, \emph{Reg. Stud.} 58, 601--618.
https://doi.org/\href{https://doi.org/10.1080/00343404.2023.2204123}{10.1080/00343404.2023.2204123}}

\end{CSLReferences}

\section{ABOUT THE STUDENT AUTHOR}\label{about-the-student-author}

Ryan Z. Johnson is an alumnus of the University of Alaska - Anchorage.
They graduated spring of 2018 with a Bachelor's of Science in
Mathematics. During their senior year, Ryan was urged by their advisor
Dr.~Mark Fitch to pursue TDA after doing presentation for a course and
connected with Dr.~Samuel Cook. Since then, they have continually worked
with Dr.~Cook outside their profession in social services. They will be
continuing their education at the University at Albany as a graduate
student in the Master's in Data Science program. This paper provided
valuable preparation and experience for the upcoming program's
concentration areas: TDA, machine learning, and statistics.




\end{document}
